\documentclass{article}
\usepackage[utf8]{inputenc}
\usepackage{amsmath}
\usepackage{geometry}
\geometry{margin=2cm}
\begin{document}

\section*{Resolução da Questão ENEM - Ampliação de Microscópio para Visualização de Hemácias}

\subsection*{Interpretação + Estratégia + Análise}

A questão aborda a capacidade de visualização de objetos microscópicos, utilizando lentes e ampliações. Para determinar a lente objetiva adequada, é necessário calcular a ampliação mínima requerida para visualizar hemácias de diâmetro 0,007 mm, considerando o limite visual de 100 micrômetros.

\subsection*{Conteúdo + Mini Revisão}

Sabemos que a ampliação total de um sistema óptico é o produto das ampliações da ocular e da objetiva:
\[
\text{Ampliação total} = \text{ampliação ocular} \times \text{ampliação objetiva}
\]

O diâmetro da hemácia convertido para micrômetros é:
\[
0,007 \text{mm} \times 1000 = 7 \text{ micrômetros}
\]

A ampliação mínima necessária para que a hemácia seja visível, dado que o limite de visualização é 100 micrômetros, é:
\[
\text{Ampliação mínima} \approx \frac{100 \text{ micrômetros}}{7 \text{ micrômetros}} \approx 14,29 \text{ vezes}
\]

A partir disso, verificamos as ampliações finais possíveis ao combinar cada lente objetiva com a ocular de 10x:

\begin{itemize}
\item Lente I (2x): 10 x 2 = 20x
\item Lente II (10x): 10 x 10 = 100x
\item Lente III (15x): 10 x 15 = 150x
\item Lente IV (1,1x): 10 x 1,1 = 11x
\item Lente V (1,4x): 10 x 1,4 = 14x
\end{itemize}

Observa-se que a menor lente objetiva cuja ampliação total é maior ou igual a aproximadamente 14,29 é a lente com 15x (lente III). Ela proporciona uma ampliação total de 150x, que é suficiente para atender ao limite de visualização.

\subsection*{Resolução + Pulo do Gato}

A solução correta é selecionar a lente III, pois ela fornece a menor ampliação que ainda garante a visualização adequada da hemácia. Mesmo que a lente V (14x ao total) seja próxima do mínimo, ela é ligeiramente insuficiente para garantir a visualização sem esforço adicional, considerando as limitações do limite de visualização.

\subsection*{Armadores}

Potenciais armadilhas incluem:
- Subestimar a ampliação necessária, pensando que ampliações menores são suficientes.
- Esquecer de converter unidades e comparar a ampliação final com a requerida.
- Confundir a adição das ampliações, ao invés de multiplicá-las.
- Curiosidade: a compreensão de que a ampliação é o produto das ampliações da ocular e da objetiva é fundamental para essa resolução.

\end{document}